\documentclass[10pt,a4paper]{article}
\usepackage{fontspec,amsmath,amssymb,graphicx}
\usepackage[margin=18mm]{geometry}
\setmainfont{Times New Roman}
\setlength{\parindent}{0pt}
\setlength{\parskip}{5pt}
\allowdisplaybreaks
\sloppy
\emergencystretch=2em
\begin{document}
\section*{BT\_4\_4\_3\_06 — Lời giải bổ sung}
\subsection*{Kiến thức và trình tự giải}
\textbf{Bước 1 — Khử điều kiện biên không thuần nhất.} Viết $u=b+w$ với $b$ thỏa đúng biên. Dirichlet hai đầu: $b=\mu+(\nu-\mu)x/L$. Neumann hai đầu: $b=\mu x+(\nu-\mu)x^2/(2L)$. Neumann–Dirichlet: $b=\nu+\mu(x-L)$; Dirichlet–Neumann: $b=\mu+x\nu$. Robin trái–Dirichlet phải: $b=\mu+d+dx$, $d=(\nu-\mu)/(L+1)$. Neumann trái–Robin phải: $b=\nu-\mu(L+1)+\mu x$. Hai Robin: $b=\mu+d+dx$, $d=(\nu-\mu)/(L+2)$.

\textbf{Bước 2 — Tính lại nguồn và dữ liệu đầu.} Từ $u_t=Ku_{xx}+F$ suy $w_t=Kw_{xx}+Q$, $Q=F+Kb_{xx}-b_t$, $w(x,0)=f(x)-b(x,0)$. Tính $b_t,b_{xx}$ trước khi chiếu; với hàm nâng bậc hai Neumann, $Kb_{xx}$ không bằng không. Nếu dùng trạng thái dừng $P$, giải $KP''+F=0$ và đúng biên rồi đặt $w=u-P$.

\textbf{Bước 3 — Chiếu và giải ODE bằng nhân tử.} Với $N_n=\int X_n^2$, đặt $d_n=N_n^{-1}\int(f-b_0)X_n$, $Q_n(t)=N_n^{-1}\int QX_n$. Phương trình mode là $T_n'+K\lambda_nT_n=Q_n$. Nhân với $e^{K\lambda_nt}$, tích phân từ $0$ đến $t$, rồi chia:


\[
T_n=d_ne^{-K\lambda_nt}+\int_0^te^{-K\lambda_n(t-s)}Q_n(s)ds,\qquad
u=b+\sum_nT_nX_n.
\]


Đây cũng là công thức Green $u=b+\int G(f-b_0)+\int_0^t\int GQ$ cho biên đã được khử.

\textbf{Bước 4 — Thực hiện tích phân thời gian.} Với $Q_n(s)=C_ne^{-\rho s}$, đặt $d=K\lambda_n$:


\[
\int_0^te^{-d(t-s)}C_ne^{-\rho s}ds
=C_ne^{-dt}\int_0^te^{(d-\rho)s}ds
=C_n\frac{e^{-\rho t}-e^{-dt}}{d-\rho}.
\]


Nếu $d=\rho$, kết quả là $C_nte^{-dt}$. Với nguồn hằng, kết quả $C_n(1-e^{-dt})/d$ khi $d>0$; mode $d=0$ phải dùng $\int_0^tQ_0(s)ds$. Nguồn đa thức thời gian tích phân từng phần hoặc giải ODE với nghiệm riêng đa thức.

\textbf{Bước 5 — Kiểm tra và tìm giới hạn.} Mỗi mode đúng PDE, biên; $T_n(0)=d_n$ trả lại dữ liệu đầu sau khi cộng $b_0$. Tách mode hằng Neumann: $\frac d{dt}\int u=K(\nu-\mu)+\int F$. Trạng thái dừng chỉ tồn tại nếu tổng nguồn và thông lượng cân bằng. Với các mode dương, kiểm tra hội tụ đủ để lấy giới hạn đều; không bỏ một mode bằng không hoặc mode tăng.
\subsection*{Lời giải áp dụng cho bài}
\textbf{Bước 1. Giải hệ nâng hai biên Robin.}
Với $b=a+dx$, $a-d=\mu$, $a+dL+d=\nu$. Trừ hai phương trình được $d=(\nu-\mu)/(L+2)$, $a=\mu+d$. Đặt $w=u-b$, $Q=F-b_t$, $q=f-b_0$. Hàm riêng là $X=\cos kx+\sin kx/k$; biên trái đúng sẵn, biên phải cho
\[(1-k^2)\sin kL+2k\cos kL=0.\]
Lấy $k>0$. Ở $k=0$, $X=A+Bx$ cho $A=B$ và $A+B(L+1)=0$, suy $A=B=0$, nên không có trị riêng $0$.
\textbf{Bước 2. Chuẩn và Green.}
\[N(k)=\frac L2(1+k^{-2})+\frac{\sin2kL}{4k}(1-k^{-2})+\frac{\sin^2kL}{k^2},\quad
G=\sum\frac{X_n(x)X_n(y)}{N(k_n)}e^{-Kk_n^2t}.\]
Vậy $u=b+\int Gq+\int_0^t\int GQ$. Công thức trên dữ liệu gốc có hai hạng biên cùng dấu dương:
\[K\int_0^t\{G(x,0,t-s)\mu(s)+G(x,L,t-s)\nu(s)\}ds.\]
Thật vậy $G_y(0)=G(0)$ và $G_y(L)=-G(L)$ trong $K[Gu_y-G_yu]_0^L$.
\textbf{Bước 3. Tính các tích phân của câu cụ thể $K=2,L=1$, dữ liệu đầu và biên $0$.}
Tích–tổng trên $[0,1]$ cho
\[A(k)=\frac{\sin k(k^2-2\pi^2)}{k(k^2-4\pi^2)},\qquad
B(k)=\frac{(1-\cos k)(k^2-2\pi^2)}{k(k^2-4\pi^2)}.\]
Chúng lần lượt là tích phân của $\cos^2\pi y$ nhân với $\cos ky$ và $\sin ky$. Do đó $C_n=(A(k_n)+B(k_n)/k_n)/N(k_n)$, không còn tích phân không gian chưa tính.
\textbf{Bước 4. Mode nguồn.}
\[T_n=C_ne^{-2k_n^2t}\int_0^te^{(2k_n^2-1)s}ds
=C_n\frac{e^{-t}-e^{-2k_n^2t}}{2k_n^2-1},\qquad
u=\sum_{n\ge1}T_nX_n.\]
Thương bằng $te^{-t}$ khi $2k_n^2=1$. ODE và từng hàm riêng kiểm tra đầy đủ PDE, biên, mặt đầu; các mode giảm cho giới hạn $0$.\end{document}